\documentclass{article}

\usepackage{xcolor}
\usepackage{makeidx}
\makeindex
\usepackage{listings}

\title{BookML test: \texttt{lstlisting} user guide}

\author{Copyright 1996--2004, Carsten Heinz%
   \\ Copyright 2006--2007, Brooks Moses
   \\ Copyright 2013--, Jobst Hoffmann}

\usepackage{bookml/bookml}
\usepackage{framed}
\AtBeginDocument{\begin{framed}
  \begin{description}
    \item[BookML] \BookMLversion
    \item[\LaTeXML] \LaTeXMLversion{} (\csname bml@generator@engine\endcsname)
    \item[\LaTeX] \fmtversion
  \end{description}
\end{framed}}


\begin{document}

\refstepcounter{section}
\refstepcounter{subsection}
\subsection{Typesetting listings}

\lstset{language=Pascal,numbers=left,frame=tblr}

{\lstset{frame=lines}
\lstinline!var i:integer;!, \lstinline$var i:integer;$

\begin{lstlisting}
for i:=maxint to 0 do
begin
    { do nothing }
end;

Write('Case insensitive ');
WritE('Pascal keywords.');
\end{lstlisting}}

\begin{lstlisting}[numbers=none]
for i:=maxint to 0 do
begin
    { do nothing }
end;

Write('Case insensitive ');
WritE('Pascal keywords.');
\end{lstlisting}

\begin{lstlisting}[numbers=right]
for i:=maxint to 0 do
begin
    { do nothing }
end;

Write('Case insensitive ');
WritE('Pascal keywords.');
\end{lstlisting}

{\lstset{frame=trbl,framesep=0pt}
\begin{lstlisting}[firstline=2,
                    lastline=5]
for i:=maxint to 0 do
begin
    { do nothing }
end;

Write('Case insensitive ');
WritE('Pascal keywords.');
\end{lstlisting}}

{\lstset{frame=lines,framesep=0pt}
\begin{lstlisting}[firstline=2,
                    lastline=5]
for i:=maxint to 0 do
begin
    { do nothing }
end;

Write('Case insensitive ');
WritE('Pascal keywords.');
\end{lstlisting}}

{\lstset{frame=topline,framesep=0pt}
\begin{lstlisting}[firstline=2,
                    lastline=5]
for i:=maxint to 0 do
begin
    { do nothing }
end;

Write('Case insensitive ');
WritE('Pascal keywords.');
\end{lstlisting}}

{\lstset{frame=bottomline,framesep=0pt}
\begin{lstlisting}[firstline=2,
                    lastline=5]
for i:=maxint to 0 do
begin
    { do nothing }
end;

Write('Case insensitive ');
WritE('Pascal keywords.');
\end{lstlisting}}

\lstinputlisting[lastline=4,language=TeX]{listings.sty}

\subsection{Figure out the appearance}
{\lstset{% general command to set parameter(s)
       basicstyle=\small,          % print whole listing small
       keywordstyle=\color{black}\bfseries\underbar,
                                   % underlined bold black keywords
       identifierstyle=,           % nothing happens
       commentstyle=\color{red}, % white comments
       stringstyle=\ttfamily,      % typewriter type for strings
       showstringspaces=false}     % no special string spaces
\begin{lstlisting}
for i:=maxint to 0 do
begin
    { do nothing }
end;

Write('Case insensitive ');
WritE('Pascal keywords.');
\end{lstlisting}}

\subsection{Seduce to use}\label{gSeduceToUse}
{\lstset{numbers=left, numberstyle=\tiny, stepnumber=2, numbersep=5pt}
\begin{lstlisting}
for i:=maxint to 0 do
begin
    { do nothing }
end;

Write('Case insensitive ');
WritE('Pascal keywords.');
\end{lstlisting}}

{\lstset{frame=tb}
\begin{lstlisting}[float,caption=A floating example]
for i:=maxint to 0 do
begin
    { do nothing }
end;

Write('Case insensitive ');
WritE('Pascal keywords.');
\end{lstlisting}}

{\lstset{literate={:=}{{$\gets$}}1 {<=}{{$\leq$}}1 {>=}{{$\geq$}}1
{<>}{{$\neq$}}1}
\begin{lstlisting}
if (i<=0) then i := 1;
if (i>=0) then i := 0;
if (i<>0) then i := 0;
\end{lstlisting}}

\refstepcounter{section}

\setcounter{subsection}{4}

\subsection{Special characters}\label{uSpecialCharacters}
{\lstset{tabsize=2}
\begin{lstlisting}
123456789
	{ one tabulator }
		{ two tabs }
123		{ 123 + two tabs }
\end{lstlisting}}

{\lstset{showspaces=true,
        showtabs=true,
        tab=\rightarrowfill}
\begin{lstlisting}
    for i:=maxint to 0 do
    begin
	{ do nothing }
    end;
\end{lstlisting}}

\subsection{Line numbers}

{\lstset{numbers=left,numberstyle=\tiny,
    stepnumber=2,numbersep=5pt}
\begin{lstlisting}[firstnumber=100]
for i:=maxint to 0 do
begin
    { do nothing }
end;

\end{lstlisting}

\begin{lstlisting}[firstnumber=last]
Write('Case insensitive ');
WritE('Pascal keywords.');
\end{lstlisting}}

{\lstset{numbers=left,numberstyle=\tiny,stepnumber=2,numbersep=5pt}
\begin{lstlisting}[name=Test]
for i:=maxint to 0 do
begin
    { do nothing }
end;

\end{lstlisting}

\begin{lstlisting}[name=Test]
Write('Case insensitive ');
WritE('Pascal keywords.');
\end{lstlisting}}

{\lstset{numbers=left,numberstyle=\tiny,stepnumber=1,numbersep=5pt}
\begin{lstlisting}[name=Test,
    language={[ansi]C},
    linerange={1-4,6-7,10-14,
    17-19,21-22},
    firstnumber=1]
#include <stdio.h>
#include <stdlib.h>

int main(int argc,char* argv[]){
    /* declaring variables */
    int i;
    int limit;

    /* checking arguments */
    if ( argc > 1 ) {
        limit = atoi(argv[1]);
    } else {
        limit = 100;
    }

    /* counting lines */
    for (i = 1;i <= limit;i++) {
        printf("Line no. %3.0d\n", i);
    }

    return 0;
}

\end{lstlisting}

\begin{lstlisting}[language={},
    linerange={1-2,6-7},
    consecutivenumbers=false]
Line no.   1
Line no.   2
Line no.   3
Line no.   4
Line no.   5
Line no.   6
Line no.   7
\end{lstlisting}}

\subsection{Layout elements}

\begin{lstlisting}[frame=single]
for i:=maxint to 0 do
begin
    { do nothing }
end;
\end{lstlisting}

\begin{lstlisting}[frame=trBL]
for i:=maxint to 0 do
begin
    { do nothing }
end;
\end{lstlisting}

{\lstset{frameround=fttt}
\begin{lstlisting}[frame=trBL]
for i:=maxint to 0 do
begin
    { do nothing }
end;
\end{lstlisting}}

{\lstset{xleftmargin=.05\linewidth}
\begin{lstlisting}[caption={Useless code},label=useless]
for i:=maxint to 0 do
begin
    { do nothing }
end;
\end{lstlisting}}

{\lstset{xleftmargin=.05\linewidth}
\begin{lstlisting}[title={`Caption' without label}]
for i:=maxint to 0 do
begin
    { do nothing }
end;
\end{lstlisting}}

{\lstset{backgroundcolor=\color{yellow}}
\begin{lstlisting}[frame=single,
                      framerule=0pt]
for i:=maxint to 0 do
begin
    j:=square(root(i));
end;
\end{lstlisting}}

\subsection{Emphasize identifiers}

{\lstset{emph={square,root},emphstyle=\underbar}
\begin{lstlisting}
for i:=maxint to 0 do
begin
    j:=square(root(i));
end;
\end{lstlisting}}

{\lstset{emph={square},      emphstyle=\color{red},
           emph={[2]root,base},emphstyle={[2]\color{blue}}}
\begin{lstlisting}
for i:=maxint to 0 do
begin
    j:=square(root(i));
end;
\end{lstlisting}}

\subsection{Indexing}

{\lstset{index={square},index={[2]root}}
\begin{lstlisting}
for i:=maxint to 0 do
begin
    j:=square(root(i));
end;
\end{lstlisting}}

\end{document}
